\section{Machine verification in Lean 4}\label{sec:lean}

The theorems of this paper are statements of algebraic geometry and are not
formalised. What is formalised is the arithmetic on which each of them turns,
and in every case that arithmetic is finite and integral, so it can be settled
by kernel computation. The file \texttt{HodgeObstruction.lean} accompanies the
paper. It compiles against a bare Lean~4 toolchain, has no \texttt{sorry}, and
depends on no library, Mathlib included; the report at its foot states, for
each of its eighty theorems, that it depends on no axiom whatever, or on
propositional extensionality alone, which is what \texttt{decide} uses when it
reduces an equality of Booleans, or when a bound is proved from the lemmas
on natural numbers in the core library.\footnote{Sixty-nine of the eighty
report \texttt{does not depend on any axioms}; the remaining eleven report
\texttt{[propext]}. No theorem reaches for anything else, and none mentions
\texttt{sorryAx}. The workflow in \texttt{.github/workflows/lean.yml} fails the
build if either changes.} The package also builds with \texttt{lake build}
under the \texttt{lakefile.toml} distributed with it, with the same report;
each run takes about two minutes. The results of \Cref{ssec:rankzero}, the annihilator
dimension and the rank-zero obstruction, are not among the eighty: the
lemma is proved by hand, and its dimension count is verified in exact
arithmetic modulo a prime by item (XIV) of \Cref{sec:verification}.

\subsection{What is checked}

\begingroup
\small
\renewcommand{\arraystretch}{1.06}
\begin{longtable}{@{}>{\raggedright\arraybackslash}p{0.31\linewidth}>{\raggedright\arraybackslash}p{0.62\linewidth}@{}}
\caption{The eighty machine-checked theorems.}\label{tab:lean}\\
\toprule
\textbf{theorem} & \textbf{statement} \\
\midrule
\endfirsthead
\multicolumn{2}{@{}l}{\footnotesize\itshape Table \thetable, continued.}\\[2pt]
\toprule
\textbf{theorem} & \textbf{statement} \\
\midrule
\endhead
\midrule
\multicolumn{2}{r@{}}{\footnotesize\itshape continued on the next page}\\
\endfoot
\bottomrule
\endlastfoot
\texttt{qnorm\_mul\_check} & the norm on $\ZZ[\sqrt{-d}]$ is multiplicative on the data used \\
\texttt{qmul\_conj} & $\tau\bbar{\tau}$ equals the norm, as an element of $\ZZ[\sqrt{-d}]$ \\
\texttt{characters\_differ} & $\tau^{2n}\ne\Nm(\tau)^{n}$ at $\tau=2+\sqrt{-d}$, for nine fields and $n\le8$ \\
\texttt{characters\_differ'} & the same at five further witnesses \\
\texttt{bad\_witness} & the witness $\tau=1+\sqrt{-d}$ fails, at $(d,n)=(1,4)$ and $(3,3)$ \\
\texttt{bad\_witness\_exceptions} & and fails exactly when $d=1,\,4\mid n$ or $d=3,\,3\mid n$ \\
\texttt{ratio\_not\_root\_of\_unity} & $\tau^{n}\ne\bbar{\tau}^{\,n}$ at $\tau=2+\sqrt{-d}$ \\
\texttt{sign\_uniform\_and\_closed} & the sign of \Cref{thm:fmtheta} is
  uniform over index sets and equals $(-1)^{g(g+1)/2+k}$, for $g\le5$ \\
\texttt{subsets\_count} & the subset enumeration used there has $\binom{g}{k}$ members \\
\texttt{semiregularity\_target} & $\sum_{q}\binom{2n}{q}\binom{2n}{q+2}=\binom{4n}{2n-2}$ for $n\le7$ \\
\texttt{secant\_count\_}\allowbreak\texttt{below\_thresholds} &
  $2n^{2}+2n<2^{2n}<3^{2n}$ for $2\le n\le8$ \\
\texttt{thresholds\_not\_necessary} & $3<2^{2}$ and $5<3^{2}$, the two retractions of \Cref{sec:numerology} \\
\texttt{grading\_positions\_differ} & $(2n,0)$, $(0,2n)$ and $(n,n)$ are distinct cells for $n\ge1$ \\
\texttt{weil\_delta\_parity} & $\sqrt{-d}^{\,k}$ is rational exactly when $k$ is even, for $k\le20$ \\
\texttt{weil\_delta\_closed\_even} & $\sqrt{-d}^{\,2m}=(-d)^{m}$ \\
\texttt{weil\_delta\_closed\_odd} & $\sqrt{-d}^{\,2m+1}=(-d)^{m}\sqrt{-d}$ \\
\texttt{weil\_term\_counts} & each of $\omega_{1},\omega_{2}$ has $2^{2n-1}$ monomials \\
\texttt{supports\_disjoint} & no monomial of the Weil line occurs in a power of $\eta$ \\
\texttt{hodge\_dim\_count} & $\dim\Hdg^{k}=3$ for $k=n$ and $1$ otherwise, for $n\le6$ \\
\texttt{weil\_line\_hodge\_type} & the Weil line is of type $(n,n)$ exactly at balanced signature \\
\texttt{clifford\_dimensions} & $\dim C=2^{b}$, $\dim C^{+}=2^{b-1}$, $\dim\mathrm{KS}=2^{b-2}$ \\
\texttt{k3\_type\_only\_at\_one} & $h^{2,0}=1$ on the pieces of $H^{2}$ forces $n=1$ \\
\texttt{moduli\_codimension} & the Weil family has codimension $n(n+1)$ in $\cA_{2n}$ \\
\texttt{discriminant\_is\_a\_norm} & $d^{n}p^{2}$ is a norm from $K$, with an explicit witness \\
\texttt{split\_codimension} & the split locus has codimension $n(n-1)/2$, zero only at $n=1$ \\
\texttt{quat\_class\_odd} & for $n$ odd the quaternionic class is $b$ times a square \\
\texttt{quat\_class\_even} & for $n$ even it is a square \\
\texttt{quat\_product\_class} & $d^{\,n-1}$ is a norm for $n$ even, with a witness \\
\texttt{quat\_locus\_codimension} & the quaternionic locus has dimension $n(n+1)/2$, codimension $n(n-1)/2$ \\
\texttt{quat\_locus\_}\allowbreak\texttt{divisor\_at\_two} & at $n=2$ that locus is a divisor: the three numbers are $4$, $3$, $1$ \\
\texttt{irrational\_}\allowbreak\texttt{witness\_exists} & some $m\le7$ makes $(m+\sqrt{-d})^{2n}$ irrational, for $n\le8$ \\
\texttt{weil\_line\_spanned} & the determinant of $w,\tau w$ is $b(w_{1}^{2}+dw_{2}^{2})$, hence nonzero \\
\texttt{norm\_form\_positive} & $a^{2}+db^{2}>0$ for every nonzero $(a,b)$, for ten fields \\
\texttt{norm\_sum\_}\allowbreak\texttt{forces\_zero} & a positive combination of norms vanishes only when every term does \\
\texttt{semiregularity\_}\allowbreak\texttt{source\_target} & at $n=2$ the source of $\sigma$ overtakes its target at $s=5$ \\
\texttt{tensor\_rank\_gap} & $\binom{2n}{2}\ge6>1$ for $n\ge2$, the separation of \Cref{lem:tensorrank} \\
\texttt{secant\_witness\_}\allowbreak\texttt{at\_n\_four} & the line bundle witness of \Cref{thm:secantexists} reproduces $6(u_{t}+v_{t})$ in every degree \\
\texttt{secant\_witness\_rank} & that witness has rank $6=Ma$, which is positive, so \Cref{lem:posrank} applies \\
\texttt{vandermonde\_}\allowbreak\texttt{nodes\_distinct} & the Vandermonde product on the nodes $0,\dots,n$ is nonzero for $n\le8$ \\
\texttt{level\_sums\_forced} & the level sums of \Cref{thm:pterange} satisfy $\sum_{j}M_{j}N_{j}^{r}=0$, $r=1,2,3$, at six level sets \\
\texttt{level\_sums\_totals} & their totals are $1$, $3$ and $-2$, so none of those objects has rank zero \\
\texttt{level\_matrix\_}\allowbreak\texttt{nonsingular} & with at most $n$ distinct norms the matrix $(N_{j}^{r})$ is nonsingular \\
\texttt{gauss\_norm\_counts} & the number of Gaussian integers of norm $1,2,3,4,5,9,45$ \\
\texttt{norm\_supply\_}\allowbreak\texttt{blocks\_small\_case} & at the nodes $1,2,3,4$ the forced $|M_{3}|=4$ while $\ZZ[i]$ has no element of norm $3$ \\
\texttt{pencil\_index\_d1} & at $d=1$ two minors of the pencil differ by the constant $2$ for $|k|\le200$, so the index of $\ZZ[i]x_{k}$ divides $N_{1}=2$ \\
\texttt{pencil\_index\_d3} & at $d=3$ a constant minor $-4$ and a second minor have gcd dividing $4$ for $|k|\le200$, so $N_{3}=4$ \\
\texttt{pencil\_beta\_positive} & $4\beta_{1}(k)=2+2k+k^{2}$ and $9\beta_{3}(k)=1+2k^{2}$ are positive for $|k|\le200$, with discriminants $-4$ and $-8$ \\
\texttt{pencil\_square\_members} & $\beta_{1}(-1)=1/4$ and $\beta_{3}(-2)=1$, the members of the pencil with a rational point \\
\texttt{split\_obstruction\_}\allowbreak\texttt{rank} & $n^{2}-n(n+1)/2=n(n-1)/2$ and $n(n-1)=2\cdot n(n-1)/2$ for $n\le30$: the rank of the obstruction map of \Cref{prop:splitobstruction}(ii) and the even side of \Cref{cor:pteobstructed} \\
\texttt{split\_obstruction\_}\allowbreak\texttt{count} & the bookkeeping of \Cref{prop:splitobstruction}(iv) at $n=3$: kernel $15=9+6$, image $3$, direct sum $12$, complement $3$ \\
\texttt{evaluation\_not\_}\allowbreak\texttt{surjective} & $s\binom{2n}{2}>2\binom{2n}{2}+4n^{2}$ for $5\le s\le40$ and $2\le n\le30$, so the evaluation map of \Cref{prop:evaluation}(iii) is not surjective; at $n=3$, $s=6$ the dimensions are $66$ and $90$ \\
\texttt{markman\_candidate\_}\allowbreak\texttt{identities} & for odd $d\le201$ and $N=(d+9)/2$: the point count $12N(N-5)=(d+9)(3d-3)$, the Euler characteristic $(d+9)(27-d)$, the two identities that place $\ch(\cF)$ at $(1,3)$, $\Nm(3+\sqrt{-d})=2N$, and the nonvanishing of the ranks, \Cref{prop:markman8}\\
\texttt{secant\_kernel\_}\allowbreak\texttt{arithmetic} & the matrix $\bigl(\begin{smallmatrix}a&b\\ b&-ad\end{smallmatrix}\bigr)$ has determinant $-(a^{2}d+b^{2})\ne0$ for $(a,b)\ne(0,0)$, $a,b<13$, $d\le20$; the coefficients $C_{k}=k!\,c_{k}$ of $au_{t}+bv_{t}$ satisfy $C_{k+1}=-d\,C_{k-1}$; and $n(n+1)/2+n(n-1)/2=n^{2}$ for $n<60$, \Cref{thm:factor}\\
\texttt{candidate\_unnormalised\_}\allowbreak\texttt{points} & for odd $d\le201$: $(d+9)(3d-3)-(d+9)(2d-1)=(d+9)(d-2)\ge12$, so an unnormalised double point exists; normalising all of them gives $\chi=50N$, and $50N\ne72N-4N^{2}$ for every integer $N>0$, \Cref{thm:candidatefails}, \Cref{rem:normalisedpoints}\\
\texttt{mumford\_rigidity\_}\allowbreak\texttt{sos} & a positive multiple of each of the seventeen forms $P_{1},P_{2},Q_{1},\dots,Q_{4},R,S_{1},\dots,S_{4},Z_{1},\dots,Z_{6}$ is a sum of squares of integral linear forms with positive coefficients, checked on $\{0,1,2\}^{4}$, which suffices for polynomials of degree at most two in each variable, \Cref{thm:mumfordrigid} \\
\texttt{smooth\_invariants} & Noether and the self-intersection formula for a smooth support, $K^{2}+e=12\chi$ and $K^{2}-e=6(b^{2}+d)^{2}$ \\
\texttt{smooth\_bmy\_defect} & the Bogomolov--Miyaoka--Yau defect of a smooth support is $24(b^{4}-d^{2})$ \\
\texttt{smooth\_bmy\_}\allowbreak\texttt{is\_d\_le\_bsq} & so the inequality is exactly $d\le b^{2}$, over a box in $b$ and $d$ \\
\texttt{smooth\_index\_vacuous} & the Hodge index bound $9(b^{2}+d)\ge8b^{2}$ holds for every $d\ge1$ \\
\texttt{smooth\_four\_}\allowbreak\texttt{discriminants} & the squarefree odd $d\le9$ are exactly $1,3,5,7$, which is \Cref{cor:fourdiscriminants} \\
\texttt{smooth\_table} & the invariants of \Cref{tab:smoothsupport} at $b=3$ \\
\texttt{burch\_weight\_}\allowbreak\texttt{identities} & the weight $x^{2}+d$ annihilates three moments, $c_{k+2}+dc_{k}=0$, which is \Cref{thm:noci} \\
\texttt{burch\_rank\_one\_}\allowbreak\texttt{discriminant} & nine times the discriminant of the rank one quadratic is $-2b^{2}-18d$, always negative \\
\texttt{burch\_rank\_four\_}\allowbreak\texttt{witness} & the rank four candidate at $d=23$ solves the moment system and is removed only by injectivity \\
\texttt{closure\_omits\_}\allowbreak\texttt{conjecture} & the Hodge conjecture is not in the forward closure of what this paper proves and quotes \\
\texttt{frontier\_}\allowbreak\texttt{suffices} & it is in the closure once (F3) alone is adjoined, which is \Cref{prop:f3ishc}; once $\{\mathrm{L},\mathrm{M}\}$, $\{\mathrm{L},\mathrm{F3}'\}$, $\{\mathrm{V},\mathrm{F3}'\}$ or $\{\mathrm{P2}_{\mathrm{split}},\mathrm{F2},\mathrm{F3}'\}$ is, with $\mathrm{P2}_{\mathrm{split}}$ the propagation statement for the split families of the CM fields of degree at least four; and, not minimally, once $\{\mathrm{P2}_{\mathrm{split}},\mathrm{F2},\mathrm{F3}\}$, $\{\mathrm{L},\mathrm{F3}\}$ or $\{\mathrm{V},\mathrm{F3}\}$ is \\
\texttt{frontier\_}\allowbreak\texttt{minimal} & in each of the five minimal sets of \Cref{thm:closure}(ii) every element is necessary: dropping any one leaves the conjecture underivable \\
\texttt{frontier\_}\allowbreak\texttt{smallest} & of the sixteen open statements exactly one, (F3), suffices alone, and of the one hundred and twenty pairs exactly those containing (F3) and the pairs $\{\mathrm{L},\mathrm{M}\}$, $\{\mathrm{L},\mathrm{F3}'\}$, $\{\mathrm{V},\mathrm{F3}'\}$ do \\
\texttt{variational\_gives\_}\allowbreak\texttt{abelian} & (V) alone puts the conjecture for abelian varieties in the closure, and not $\mathbf{HC}$, \Cref{thm:vhcab} \\
\texttt{lefschetz\_gives\_}\allowbreak\texttt{abelian} & (L) gives (V), hence the conjecture for abelian varieties, and not $\mathbf{HC}$, \Cref{prop:lefvhc} \\
\texttt{secant\_route\_}\allowbreak\texttt{descends} & granting both open demands of the secant route reaches $\delta=\delta_{0}$ and, through \Cref{prop:descent}, every imaginary quadratic family, and neither the CM fields of higher degree nor the conjecture \\
\texttt{propagation\_}\allowbreak\texttt{closes} & (P2) for the split imaginary quadratic families puts the whole imaginary quadratic case in the closure, and (P2) for the split families of the fields of degree at least four puts the Weil classes of every CM field there \\
\texttt{p2\_minimal\_}\allowbreak\texttt{count} & $\binom{4n}{2}-4n^{2}=2\binom{2n}{2}=2n(2n-1)$ for every $n\le60$, the dimension at which the first form of the criterion would be met, \Cref{thm:p2numerical}, and which no complex attains, \Cref{thm:p2false} \\
\texttt{mumford\_}\allowbreak\texttt{invariants} & the invariants of $\mathfrak{sl}_{2}^{3}$ in $\bigwedge^{q}(V_{1}\otimes V_{2}\otimes V_{3})$ have dimensions $1,0,1,0,1,0,1,0,1$, counted from weight multiplicities, \Cref{cor:mumfordlefschetz} \\
\texttt{two\_branch\_}\allowbreak\texttt{annihilator} & for $n=2,3$ the monomials of type $(p,p)$ killed by $Q\wedge HH^{1}$, $P\wedge HH^{1}$ and $HH^{1}\wedge HH^{1}$ are $\alpha_{-}$, $\alpha_{+}$ and none, \Cref{thm:nosum} \\
\texttt{weil\_monomials\_}\allowbreak\texttt{separated} & neither Weil monomial lies in $\bigwedge^{*}(V_{-}^{1,0}\oplus V_{+}^{0,1})$ or $\bigwedge^{*}(V_{+}^{1,0}\oplus V_{-}^{0,1})$, for $n\le12$, which is the last step of \Cref{thm:p2support}(iii) \\
\texttt{mumford\_ext\_}\allowbreak\texttt{profile} & the bounds $1,16,119,328,560,328,119,16,1$ are palindromic, with alternating sum $112$, sum $1488$ and $2+238+560=800$, \Cref{thm:mumfordprofile} \\
\texttt{lefschetz\_}\allowbreak\texttt{counts} & the $\Sp_{8}$-invariants of $\bigwedge^{*}(V\oplus V)$ number $1,3,6,10,15,10,6,3,1$ by the decomposition of each $\bigwedge^{i}V$, the Hodge classes that are not Lefschetz $0,0,2,6,13,6,2,0,0$, the Lefschetz classes at a CM point are counted by $(1+4y+y^{2})^{4}$, and Weyl's formula gives $308$ for $2\varpi_{2}$, \Cref{thm:lefschetzclosure} \\
\texttt{criterion\_}\allowbreak\texttt{endomorphisms\_n2} & for all natural numbers, $2e_{0}+12=\chi+2e_{1}$ with $e_{1}\ge8$ and $\chi\ge1$ gives $e_{0}\ge3$: the Euler characteristic of an object meeting the criterion at $n=2$, \Cref{thm:p2endo}(iv) \\
\texttt{criterion\_}\allowbreak\texttt{endomorphisms\_n3} & for all natural numbers, $2e_{0}+60=\chi+2e_{1}+e_{3}$ with $e_{1}\ge12$, $e_{3}\ge40$ and $\chi\ge1$ gives $e_{0}\ge3$: the same at $n=3$, \Cref{thm:p2endo}(iv) \\
\end{longtable}
\endgroup

An element of $\ZZ[\sqrt{-d}]$ is a pair of integers, multiplication and
conjugation are the evident operations, and every statement in the table is a
closed Boolean expression evaluated by the kernel.%
\footnote{The tactic is \texttt{decide} throughout, which asks the kernel to evaluate a closed Boolean expression by reduction. That is the strongest check available and the least flexible: it works only over finite ranges, which is why every theorem in the table carries an explicit bound, and it is also why none of them can be the general statement the paper proves.} Nothing is quoted; the
kernel recomputes each assertion from the definitions.

Two of the entries deserve a word about what they do and do not settle.
\texttt{hodge\_dim\_count} verifies the bookkeeping of
\Cref{thm:hdgcount}, namely that summing the contributions
$(n,n)$, $(2n,0)$ and $(0,2n)$ over the degrees gives three in the middle and
one elsewhere; the representation theory that identifies those contributions is
in the proof of \Cref{thm:hdgcount} and is not formalised. Likewise
\texttt{k3\_type\_only\_at\_one} verifies that the equation $h^{2,0}=1$ has
$n=1$ as its only solution once the three possible values $0$, $n^{2}$ and
$n(n-1)$ are known; the identification of those values is
\Cref{prop:normpart}. In both cases the division of labour is the same as
everywhere else in this section: the kernel settles the computation, and the
mathematics around it is written out and must be read.

\subsection{Two statements that look true and are not}

Two theorems of the file record statements that a careful reader would accept
and that are false, and both bear on \Cref{sec:characters}.

The first concerns the witness $\tau=1+\sqrt{-d}$ for \Cref{lem:charsdiffer}.
It is the obvious choice, and it does not work at every field:
\texttt{bad\_witness} shows that at $d=1$ the ratio $\tau/\bbar{\tau}$ is
$\sqrt{-1}$ and at $d=3$ it is a primitive cube root of unity, so at those
fields the two characters agree at that particular $\tau$ for suitable $n$.
The lemma is true, but not by that route, which is why the proof given in
\Cref{sec:characters} uses only the finiteness of $\mu(K)$ and avoids
witnesses altogether.

The second concerns the exceptional set of that witness at $d=3$, which is
$3\mid n$ and not $6\mid n$, as a hasty reading of the order of $\mu(K)$
would suggest; \texttt{bad\_witness\_exceptions} records the correct
condition. Neither affects \Cref{thm:character}. Both are the kind of slip that
a careful reading does not catch and a kernel check does, which is the reason
for keeping them in the file.

\subsection{What is not formalised}

Everything else, and we are explicit about it. The cohomology of an abelian
variety as an exterior algebra, the Hodge decomposition, the Fourier--Mukai
transform as a functor, Grothendieck--Riemann--Roch, the Lefschetz hyperplane
theorem, hard Lefschetz, the invariant theory of $\SL_{2n}$, the Clifford
algebra as a functor, and the Lefschetz standard conjecture are all used as
mathematics and none is formalised here. A formalisation of those would
require a library of complex algebraic geometry that does not presently exist.
What the Lean file establishes is that where this paper computes, it computes
correctly; where it reasons, the reader must still read.
